\section{Introduction}
\IEEEPARstart{C}{anteen} meals are the main source of daily energy for most
Chinese university students, yet the energy content of a dish is rarely visible
at the point of choice. A system that takes a single photograph of a tray and
returns the three most likely dishes together with a rough energy interval would
make calorie awareness immediate, requires no wearable sensor and no manual
logging, and can run on a commodity laptop or an edge device.

Building such a system is not simply an image-classification exercise, and the
difficulty is not the one a Western benchmark suggests. First, the menu is
Chinese: the dishes a canteen serves---mapo tofu, braised pork, stir-fried potato
shreds, steamed egg---are almost absent from Food-101, so a model trained on the
usual benchmark answers a question nobody asked. Second, the photographs are
taken with hand-held phones under mixed artificial light, at oblique angles, with
cluttered trays, partially eaten portions and steam, and a classifier trained on
curated studio photographs degrades sharply under this shift
\cite{hendrycks2019robustness}. Third, and most awkwardly for a deployed system,
a softmax classifier always returns a class: when the tray holds a dish outside
the menu, or no dish at all, the model still produces a confident answer. Since
the number of dishes a canteen serves is large and the frequency distribution is
skewed, the useful question is not only \emph{which} dish but \emph{whether to
answer at all}.

This paper therefore studies three things that a course-scale project can
actually control. We ask how much Chinese pre-training is worth by fitting the
same canteen head twice, once from ImageNet weights and once from a backbone
pre-trained on 241 Chinese dish classes. We ask how much test-time fusion of
several framings is worth, since a tray photograph has no canonical crop. And we
ask how a deployed system should decline to answer, which we treat as a
measurable quantity on five evaluation sets rather than as a threshold chosen by
eye.

Our contributions are:
\begin{itemize}
\item a 21-way Chinese canteen protocol built on CNFOOD-241 \cite{cnfood241}
with a leak-free split, a 241-class pre-training stage, and a catch-all bucket
that makes "some other dish" an explicit class instead of a silent error;
\item a controlled measurement of pre-training scale: the identical recipe and
head move from \cnNoPreTest\% to \cnTest\% test accuracy when the backbone is
pre-trained on Chinese dishes (\cnGainPretrain{} points), which is larger than
any adaptation-policy change we could obtain on the Western benchmark;
\item a six-view test-time fusion whose gain (\cnFuseGain{} points) is measured
against the single-view protocol used for model selection, and a per-class error
analysis that locates the residual error in intra-class variation
(sweet-and-sour ribs versus sweet-and-sour pork, twice-cooked pork versus
braised pork) rather than in class imbalance;
\item an open-set gate that combines a logistic domain probe on pooled features,
a Chinese-dish veto from a wider head, and a catch-all posterior, evaluated on
five sets (canteen validation, canteen test, fifteen unseen Chinese dishes,
non-Chinese food, non-food scenes and real photographs) with the coverage it
costs on the supported set reported alongside the rejections it buys;
\item an external evaluation on \webN{} web images of the twenty dishes,
screened for exact and near duplicates against the training data, which is the
only number in this paper produced by images that no part of the pipeline has
ever seen;
\item three negative results---stronger photometric augmentation,
class-balanced loss and focal loss---each of which we report with the same
protocol that produced the positive ones.
\end{itemize}

Section \ref{sec:data} describes the data and the splits, Section
\ref{sec:method} the pipeline, Section \ref{sec:exp} the experiments, and
Section \ref{sec:lim} the limitations we consider binding.
